\subsection{Analisis Kebutuhan Infrastruktur Audit Trail}

Selain mengidentifikasi \textit{audit event} yang perlu dicatat, perancangan sistem \textit{audit trail} juga memerlukan analisis terhadap infrastruktur yang dibutuhkan untuk mendukung proses pembangkitan, penyimpanan, pengelolaan, dan pemanfaatan \textit{audit record}. Analisis ini dilakukan dengan mengacu pada keluarga kontrol Audit and Accountability (AU) pada NIST SP 800-53 Rev. 5 \cite{nist80053r5} serta panduan manajemen log pada NIST SP 800-92 \cite{nist80092}. Berdasarkan kedua acuan tersebut, diidentifikasi sejumlah kapabilitas yang perlu dimiliki oleh sebuah sistem \textit{audit trail}. Kapabilitas tersebut kemudian digunakan untuk menentukan komponen infrastruktur yang diperlukan serta menjadi dasar dalam mengevaluasi apakah arsitektur DecMed telah memenuhi kebutuhan tersebut.

Tabel~\ref{tab:audit-infrastructure-requirement} menyajikan kebutuhan kapabilitas infrastruktur \textit{audit trail} beserta komponen yang diperlukan untuk memenuhi masing-masing kapabilitas.

\begin{xltabular}
{\textwidth}
{|>{\centering\arraybackslash}p{1.2cm}
 |>{\raggedright\arraybackslash}p{4.5cm}
 |>{\raggedright\arraybackslash}X
 |>{\centering\arraybackslash}p{2cm}|}

\caption{Kebutuhan Infrastruktur Audit Trail}
\label{tab:audit-infrastructure-requirement}\\

\hline
\textbf{ID} &
\textbf{Kapabilitas yang Dibutuhkan} &
\textbf{Komponen Infrastruktur yang Diperlukan} &
\textbf{Control AU} \\
\hline
\endfirsthead

\hline
\textbf{ID} &
\textbf{Kapabilitas yang Dibutuhkan} &
\textbf{Komponen Infrastruktur yang Diperlukan} &
\textbf{Control AU} \\
\hline
\endhead

\hline
\endfoot

\hline
\endlastfoot

IR1 &
Menghasilkan audit record dari audit event yang terjadi pada seluruh komponen sistem &
Audit Event Source &
AU-2 \\
\hline

IR2 &
Mengumpulkan audit record dari berbagai komponen secara terpusat &
Audit Collection Service &
AU-12 \\
\hline

IR3 &
Menyimpan audit record secara terpusat dan mendukung sifat append-only &
Audit Repository &
AU-9 \\
\hline

IR4 &
Menjamin integritas audit record selama masa penyimpanan &
Integrity Protection Mechanism &
AU-9 \\
\hline

IR5 &
Menerapkan kebijakan retensi terhadap audit record &
Retention Management &
AU-11 \\
\hline

IR6 &
Menyediakan mekanisme pencarian dan peninjauan audit record &
Audit Query Interface &
AU-6 \\
\hline

IR7 &
Menyediakan mekanisme ekspor audit record untuk kebutuhan audit maupun investigasi &
Audit Export Interface &
AU-6 \\
\hline

\end{xltabular}

Berdasarkan kebutuhan infrastruktur tersebut, dilakukan evaluasi terhadap arsitektur DecMed untuk mengidentifikasi komponen yang telah tersedia maupun komponen yang masih belum tersedia. Hasil evaluasi tersebut disajikan pada Tabel~\ref{tab:gap-infrastruktur}.

\begin{xltabular}
{\textwidth}
{|>{\centering\arraybackslash}p{1cm}
 |>{\raggedright\arraybackslash}p{2.5cm}
 |>{\raggedright\arraybackslash}p{4.5cm}
 |X|}

\caption{Analisis Gap Infrastruktur Audit Trail pada DecMed}
\label{tab:gap-infrastruktur}\\

\hline
\textbf{ID} &
\textbf{Kapabilitas} &
\textbf{Kondisi pada DecMed} &
\textbf{Analisis Gap} \\
\hline
\endfirsthead

\hline
\textbf{ID} &
\textbf{Kapabilitas} &
\textbf{Kondisi pada DecMed} &
\textbf{Analisis Gap} \\
\hline
\endhead

\hline
\endfoot

\hline
\endlastfoot

IR1 &
Audit Event Source &
Seluruh aktivitas sistem telah dihasilkan oleh berbagai komponen seperti Client, PRE Server, Gas Station, IPFS, dan IOTA. &
Tidak terdapat gap. Audit event telah dihasilkan oleh komponen-komponen DecMed. \\
\hline

IR2 &
Audit Collection Service &
Setiap komponen bekerja secara independen dan belum terdapat mekanisme yang mengumpulkan audit event secara terpusat. &
Diperlukan komponen audit collection yang mampu menerima audit event dari seluruh komponen DecMed. \\
\hline

IR3 &
Audit Repository &
DecMed telah memiliki IOTA On-Chain State untuk menyimpan PatientAccessLog serta IPFS sebagai penyimpanan objek off-chain. &
Tidak terdapat gap pada media penyimpanan. Repository yang tersedia dapat dimanfaatkan untuk menyimpan audit trail, namun perlu dilakukan perancangan skema penyimpanan audit record agar memenuhi kebutuhan audit trail yang telah diidentifikasi.\\
\hline

IR4 &
Integrity Protection &
Integritas transaksi dijamin oleh IOTA, namun tidak tersedia mekanisme perlindungan integritas terhadap keseluruhan audit record. &
Diperlukan mekanisme perlindungan integritas audit record secara menyeluruh. \\
\hline

IR5 &
Retention Management &
Tidak terdapat mekanisme pengelolaan retensi audit record. &
Diperlukan mekanisme retensi sesuai kebutuhan audit dan regulasi. \\
\hline

IR6 &
Audit Query Interface &
PatientAccessLog hanya dapat diakses melalui fungsi tertentu dan tidak mendukung pencarian audit secara umum. &
Belum tersedia antarmuka pencarian audit record yang terintegrasi. \\
\hline

IR7 &
Audit Export Interface &
Tidak tersedia mekanisme ekspor audit record. &
Diperlukan mekanisme ekspor audit record untuk mendukung proses audit maupun investigasi. \\
\hline

\end{xltabular}

Berdasarkan hasil analisis pada Tabel~\ref{tab:gap-infrastruktur}, diketahui bahwa DecMed telah memiliki sumber audit event yang berasal dari berbagai komponen sistem. Namun demikian, sebagian besar kapabilitas infrastruktur yang dibutuhkan oleh sistem audit trail belum tersedia, khususnya mekanisme pengumpulan audit secara terpusat, repositori audit, pengelolaan retensi, antarmuka pencarian, serta mekanisme ekspor audit record. Selain itu, mekanisme yang telah tersedia, seperti \textit{PatientAccessLog} dan riwayat transaksi pada jaringan IOTA, hanya memenuhi sebagian kebutuhan audit trail dan belum dirancang sebagai infrastruktur audit yang terintegrasi. Oleh karena itu, diperlukan perancangan infrastruktur audit trail yang mampu memenuhi seluruh kebutuhan tersebut pada tahap rancangan solusi.